\newcommand{\auditoriaid}{A08 -- Próximos pasos, conclusiones y recomendaciones}
% Preambulo comun de las auditorias del informe E1 (revision_e1/)
\documentclass[11pt,a4paper]{article}
\usepackage[utf8]{inputenc}
\usepackage[T1]{fontenc}
\usepackage[spanish,es-noshorthands]{babel}
\usepackage{lmodern}
\usepackage{microtype}
% Sin patrones de guionado en espanol en esta instalacion (ver A10): se
% desactiva el guionado automatico para evitar cortes con reglas del ingles.
\hyphenpenalty=10000
\exhyphenpenalty=10000
\usepackage[a4paper,margin=2.2cm]{geometry}
\usepackage{booktabs,array,longtable}
\usepackage{graphicx,caption,subcaption}
\usepackage{xcolor}
\usepackage{enumitem}
\usepackage{fancyhdr}
\usepackage{amsmath}
\usepackage{tikz}
\usetikzlibrary{arrows.meta,positioning}
\usepackage{natbib}
\usepackage{xurl}
\usepackage[hidelinks]{hyperref}
\urlstyle{tt}
\newcommand{\archivo}[1]{\url{#1}}
\newcolumntype{L}[1]{>{\raggedright\arraybackslash}p{#1}}
\definecolor{alta}{RGB}{180,30,30}
\definecolor{media}{RGB}{200,120,0}
\definecolor{baja}{RGB}{60,110,160}
\definecolor{cajafondo}{RGB}{245,247,250}
\newcommand{\sevA}{\textcolor{alta}{\textbf{Alta}}}
\newcommand{\sevM}{\textcolor{media}{\textbf{Media}}}
\newcommand{\sevB}{\textcolor{baja}{\textbf{Baja}}}
\setlength{\parskip}{0.35em}
\setlength{\emergencystretch}{3em}
\setlength{\parindent}{0pt}
\pagestyle{fancy}
\fancyhf{}
\fancyhead[L]{\small Auditoría del Informe del Entregable~1 -- \auditoriaid}
\fancyhead[R]{\small \thepage}
\renewcommand{\headrulewidth}{0.4pt}
% Entorno de hallazgos: ID | Severidad | Ubicacion | Hallazgo y evidencia | Correccion
\newenvironment{hallazgos}{%
\small
\begin{longtable}{@{}L{1.15cm}L{1.25cm}L{1.9cm}L{6.0cm}L{\dimexpr\textwidth-10.3cm-10.6\tabcolsep\relax}@{}}
\toprule
\textbf{ID} & \textbf{Sev.} & \textbf{Ubicación} & \textbf{Hallazgo y evidencia} & \textbf{Corrección aplicada} \\
\midrule\endfirsthead
\toprule
\textbf{ID} & \textbf{Sev.} & \textbf{Ubicación} & \textbf{Hallazgo y evidencia} & \textbf{Corrección aplicada} \\
\midrule\endhead
\bottomrule\endfoot
}{\end{longtable}}
% Macros compartidas con el informe revisado (las secciones propuestas las usan)
% Macros compartidas entre el informe revisado y las auditorias.
% Cifras verificadas contra los archivos del repositorio (ver A01/A04/A07):
\newcommand{\nUniverso}{102}   % source_id CMIP6 con tos/Omon en ESGF (00b, model_availability_priority.csv)
\newcommand{\nSeed}{41}        % publican los cuatro experimentos (config/models_seed_cmip6.csv)
\newcommand{\nSel}{40}         % completos + QC PASS (models_inventory_final.csv)
\newcommand{\nComb}{160}       % combinaciones modelo-experimento PASS (qc_report.csv)
\newcommand{\nNoEnc}{44}       % config/models_no_encontrado_44.csv
\newcommand{\nParcial}{18}     % 102 - 40 - 44
% Fecha de la portada: fija (no \today), editable en un solo lugar.
\newcommand{\fechaentrega}{28 de septiembre de 2026}
% Estado de codigo que documenta el E1 (antes de los cambios del E2)
\newcommand{\commitEone}{\texttt{8ee6184}}

% En las auditorias el texto propuesto se muestra fuera del informe: las
% referencias cruzadas se imprimen con el nombre de la seccion destino.
\makeatletter
\def\etq#1#2{\expandafter\def\csname etq@#1\endcsname{#2}}
\etq{sec:resumen}{Resumen}
\etq{sec:alcance}{Alcance}
\etq{sec:marco}{Revisión científica}
\etq{sec:antecedentes}{Antecedentes}
\etq{sec:referencias}{Literatura sobre TOE}
\etq{tab:referencias}{bibliografía TOE}
\etq{sec:similitudes}{Niño~3.4 frente a Niño~1+2}
\etq{sec:toe_metodos}{Elección de los métodos de TOE}
\etq{tab:metodos_toe}{métodos TOE}
\etq{eq:szabo_signal}{ecuación señal M1}
\etq{eq:szabo_V}{ecuación V(t)}
\etq{eq:szabo_T}{ecuación T(t)}
\etq{eq:szabo_toe}{ecuación TOE M1}
\etq{eq:tod}{ecuación ToD}
\etq{sec:datos}{Datos}
\etq{fig:cmip6_contribution}{contribución institucional}
\etq{sec:tabla_modelos}{Modelos seleccionados}
\etq{tab:modelos}{modelos}
\etq{sec:desviacion}{Selección de modelos}
\etq{sec:dominios_periodos}{Dominios y periodos}
\etq{fig:region_nino}{regiones Niño}
\etq{sec:arquitectura}{Arquitectura del pipeline}
\etq{fig:flujograma}{Flujograma}
\etq{sec:scripts}{Descripción de los pasos}
\etq{sec:paso00b}{Paso 00b}
\etq{sec:paso01}{Paso 01}
\etq{sec:paso02}{Paso 02}
\etq{sec:paso04}{Paso 04}
\etq{sec:paso05}{Paso 05}
\etq{sec:paso06}{Paso 06}
\etq{sec:paso07}{Paso 07}
\etq{sec:reporte}{Reporte de estado}
\etq{sec:resultados}{Resultados}
\etq{fig:qc}{matriz de estado}
\etq{tab:cuentas}{conteo de modelos}
\etq{fig:mapas2d}{mapas 2D}
\etq{fig:series}{series}
\etq{fig:boxplot}{boxplots}
\etq{sec:roadmap}{Próximos pasos}
\etq{tab:roadmap}{próximos pasos}
\etq{sec:conclusiones}{Conclusiones}
\etq{sec:recomendaciones}{Recomendaciones}
\makeatother

\makeatletter
\AtBeginDocument{\renewcommand{\ref}[1]{\@ifundefined{etq@#1}{\textit{[#1]}}{\textit{\csname etq@#1\endcsname}}}}
\makeatother
% Texto propuesto: secciones degradadas un nivel y en caja
\newenvironment{propuesto}{%
  \par\medskip\noindent\textcolor{baja}{\rule{\textwidth}{0.8pt}}\par
  \noindent{\small\textbf{Inicio del texto propuesto para el informe revisado}}\par\medskip
  \begingroup
  \let\section\subsection\let\subsection\subsubsection\let\subsubsection\paragraph
}{%
  \endgroup
  \par\noindent{\small\textbf{Fin del texto propuesto}}\par
  \noindent\textcolor{baja}{\rule{\textwidth}{0.8pt}}\par\medskip}
% Recuadro de actualizacion posterior (decisiones del autor)
\newcommand{\actualizacion}[1]{\par\medskip\noindent\fcolorbox{media}{media!8}{\parbox{\dimexpr\linewidth-2\fboxsep-2\fboxrule\relax}{\small\textbf{Actualización del 29-09-2026 (decisiones del autor).} #1}}\par\medskip}

\begin{document}

\begin{center}
{\Large\bfseries Auditoría A08: Próximos pasos, conclusiones y recomendaciones}\\[0.4em]
{\normalsize Informe del Entregable~1, Secciones 8--10 (\texttt{informe\_e1\_smnv3.tex}, commit \texttt{330fd2e})}\\[0.2em]
{\small Revisión: \fechaentrega}
\end{center}

\section{Objeto}

Se auditan las secciones de cierre. Las conclusiones deben apoyarse solo en lo que el informe demuestra, y la hoja de ruta debe coincidir con lo que efectivamente se está haciendo en el Entregable~2. Para lo segundo se contrastó la tabla ``Próximos pasos'' con \archivo{informe/borradores/marco_teorico_e2.txt} y con \archivo{scripts_e2/} (pasos p01--p09). Las conclusiones se revisaron contra los hallazgos de A01--A07.

\section{Contraste de la hoja de ruta con el Entregable~2 en curso}

\begin{center}\small
\begin{tabular}{@{}L{5.2cm}L{\dimexpr\textwidth-5.2cm-2\tabcolsep\relax}@{}}
\toprule
Informe E1 (fila E2) & Metodología real del E2 \\
\midrule
``corrección de sesgo (Linear Scaling)'' & \textbf{Descartada.} El E2 muestra que la LS se cancela algebraicamente en toda anomalía; los índices se calculan con la climatología propia de cada fuente y el sesgo se reporta como diagnóstico aparte. \\
``diagrama de Taylor sobre los eventos cálidos extremos'' & Taylor sobre la serie completa (p05) y sobre el compuesto del evento típico (p08), más el puntaje de habilidad de Taylor (2001) para la selección final (p09). \\
ONI/RONI (Niño~3.4), ICEN (Niño~1+2) & Coincide (p01, p07). \\
(no mencionado) & Correlación (p06), variabilidad interanual (p04), clasificación de eventos con umbral dual (p07). \\
\bottomrule
\end{tabular}
\end{center}

\section{Hallazgos}

\begin{hallazgos}
A08-1 & \sevA & Cuadro 8, fila E2 & Anuncia ``corrección de sesgo (Linear Scaling)'', que el E2 descartó con fundamento. Si el informe E1 se entrega así, contradice al E2. & Fila reescrita según la metodología real del E2. \\
A08-2 & \sevM & Cuadro 8, título & ``Gráficas y comparaciones exigidas por el TdR'': el cuadro incluye elecciones del equipo (Taylor, RONI, análisis compuesto, gradiente zonal) que el TdR no exige. & ``Análisis previstos para los Entregables~2 a~4''. \\
A08-3 & \sevM & Fila E3 & ``gradiente zonal del Pacífico'' no está en la actividad (c) del TdR. & Se reemplaza por lo que pide la actividad (c): estado medio, variabilidad y eventos extremos con ONI e ICEN. \\
A08-4 & \sevA & Conclusión (iv) & ``un \texttt{git clone} \ldots\ más \texttt{./run.sh} reconstruye el mismo estado documentado'': desde \texttt{main} ya no es cierto (A01-6, A04-5, A05-3). & \textit{Actualización 29-09}: sin mencionar commits; se indica que el repositorio es de libre acceso y que \texttt{./run.sh} reconstruye el conjunto de datos documentado. \\
A08-5 & \sevM & Conclusión (i) & Repite por cuarta vez la oración ``40 de 102 \ldots\ 160 combinaciones'' (comentario \texttt{\% REVISAR} del Resumen). & Frase corta con remisión a Resultados. \\
A08-6 & \sevM & Conclusión (ii) & Sintaxis rota (``profundiza: la bibliografía revisada permite justificar, con evidencia explícita lo requerido por el TdR, la adopción \ldots''); llama ``prescrito'' a M1, que el propio informe presenta como adoptado; ``cuatro líneas independientes convergen'' está sobredimensionado (A03-11). & Redacción nueva; ``hipótesis de trabajo'' apoyada además en el dato propio de ERSSTv5. \\
A08-7 & \sevM & Concl. & Falta la conclusión cuantitativa que sí producen los datos del E1: el sesgo cálido de Niño~1+2 en 38 de 40 modelos (A07). & Nueva conclusión (iii), presentada como preliminar. \\
A08-8 & \sevM & Recom. & No hay recomendación sobre dos decisiones con costo documentado: exigir \texttt{ssp370} (49 → 41, A04-4) y preservar la dispersión estructural (Lois). & Dos recomendaciones nuevas. \\
A08-9 & \sevB & Título & ``Proximos pasos'' sin tilde. & ``Próximos pasos''. \\
A08-10 & \sevB & Cierre de Próximos pasos & ``El Entregable~1 queda completo y verificable sin depender de ningún paso posterior'' repite la conclusión. & Se conserva solo la frase sobre \texttt{data/processed/}. \\
\end{hallazgos}

\section{Extensión}

Original: 629 palabras (216 + 273 + 140). Propuesto: 654 
palabras. El aumento corresponde a la conclusión (iii) y a las dos recomendaciones nuevas.

\section{Texto propuesto}

\begin{propuesto}
% ============================================================
\section{Próximos pasos}
\label{sec:roadmap}
% ============================================================

El Cuadro~\ref{tab:roadmap} resume, para dar continuidad al proyecto, los análisis previstos en los Entregables~2 a~4 según las actividades~(b) a~(e) del TdR.

\begin{table}[htbp]
\caption{Análisis previstos para los Entregables~2 a~4.}
\label{tab:roadmap}
\small
\begin{tabular}{@{}L{2.2cm}L{\dimexpr\textwidth-2.2cm-2\tabcolsep\relax}@{}}
\toprule
\textbf{Entregable} & \textbf{Análisis previstos} \\
\midrule
E2 & Climatología y ciclo anual observados y simulados; sesgo, RMSE, correlación y diagrama de Taylor en ambas cajas; variabilidad interanual; eventos cálidos extremos del ENOS con ONI y RONI (Niño~3.4) e ICEN (Niño~1+2), con análisis compuesto del evento típico; puntaje de habilidad para seleccionar los modelos que mejor representan el periodo histórico. \\
E3 & Proyección del estado medio y de la variabilidad de la TSM en Niño~3.4 y Niño~1+2 bajo SSP2-4.5 y SSP5-8.5 (con SSP3-7.0 como escenario intermedio), para 2036--2065 y 2071--2100; cambios en la ocurrencia de eventos extremos del ENOS mediante ONI e ICEN. \\
E4 & Partición de la incertidumbre (variabilidad interna, modelo y escenario) y razón señal-ruido; TOE por los métodos M1 y ToD en ambas cajas; estudio técnico final con documentación metodológica, bases de datos, \emph{scripts}, conclusiones y recomendaciones, incluida la contrastación de la hipótesis de la Sección~\ref{sec:similitudes}. \\
\bottomrule
\end{tabular}
\end{table}

El Entregable~2 puede comenzar leyendo solo \texttt{data/processed/}, sin volver a ejecutar ni conocer el detalle interno de este módulo.

% ============================================================
\section{Conclusiones}
\label{sec:conclusiones}
% ============================================================

El Entregable~1, en cumplimiento de la actividad~(a) del TdR, deja construido, verificado y versionado el insumo del que dependen los Entregables~2 a~4. De su ejecución se desprenden las siguientes conclusiones:

\begin{enumerate}[label=(\roman*)]
\item El \emph{pipeline} está completo y verificado: 40 modelos CMIP6 con sus cuatro experimentos aprobados (Sección~\ref{sec:resultados}). Sus controles por valor numérico detectan problemas de los datos de origen que un control basado solo en metadatos no detectaría, como la contaminación numérica de FGOALS-f3-L y las unidades mal declaradas de GISS-E2-1-G.

\item La revisión científica (Sección~\ref{sec:marco}) sustenta la adopción de dos métodos de TOE conceptualmente independientes, M1 \citep{szabo2016} y ToD \citep{gopika2024}, y fija sus ecuaciones antes de su aplicación en el Entregable~4. De esa revisión surge además una hipótesis de trabajo contrastable: Niño~1+2 emergerá más tarde, o con mayor incertidumbre, que Niño~3.4. La mayor variabilidad interanual observada en Niño~1+2 es coherente con ella.

\item La comparación preliminar con ERSSTv5 muestra un sesgo cálido en Niño~1+2 en 38 de los 40 modelos (mediana de $+2{,}2$\,°C), mientras que en Niño~3.4 los modelos se distribuyen alrededor de la observación. La evaluación formal del Entregable~2 debe tratar ambas cajas por separado.

\item Ampliar el universo de candidatos más allá de la selección de \citet{bruno2023} deja al Entregable~2 con más información: la decisión sobre qué modelos representan mejor el ENOS se tomará con evaluación propia y no con una lista pensada para otras variables.

\item El trabajo es reproducible: el repositorio \archivo{smn_toe} es de libre acceso y, con el entorno instalado, \texttt{./run.sh} reconstruye el conjunto de datos documentado; cada corrida deja además un reporte de estado verificable. El procedimiento se detalla en el documento de anexos.
\end{enumerate}

% ============================================================
\section{Recomendaciones}
\label{sec:recomendaciones}
% ============================================================

Para la continuidad del proyecto se recomienda:

\begin{itemize}
\item Evaluar en el Entregable~2 la representación del ENOS en los 40 modelos (climatología, sesgos, métricas estadísticas e índices ONI, RONI e ICEN en ambas cajas) antes de dar por buena su idoneidad. El control de calidad de este entregable verifica la integridad y homogeneidad del dato, no su desempeño climático.
\item Verificar que la selección final de modelos del Entregable~2 conserve la dispersión estructural del ensamble \citep{lois2026}, para no subestimar la incertidumbre de modelo ni el rango del TOE.
\item Considerar la incorporación de los 8 modelos que publican \texttt{historical}, \texttt{ssp245} y \texttt{ssp585} pero no \texttt{ssp370} en los análisis que usan solo SSP2-4.5 y SSP5-8.5, lo que ampliaría los candidatos de 41 a 49 modelos.
\item Calcular y reportar M1 y ToD siempre en paralelo, documentando en el Entregable~4 los casos en que difieran sustancialmente: la bibliografía revisada muestra que esa discrepancia entre métodos es la norma y no la excepción.
\end{itemize}

\end{propuesto}

\bibliographystyle{plainnat}
\bibliography{../informe/references_rev}
\end{document}
